\section{本章小结}
\label{sec:clt-summary}

本章从掷硬币的正面次数出发，问许多独立小项之和的波动服从什么规律。第一步是选对中心和尺度：部分和的典型偏差按$\sigma\sqrt{n}$增长，平均的典型偏差按$\sigma/\sqrt{n}$缩小，只有放大$\sqrt{n}$倍的平均偏差才有不退化也不发散的极限。大数律确定了中心，中心极限定理刻画的是放大以后的形状；标准化只把均值和方差固定为$0$和$1$，形状接近钟形是独立求和的结果，特制硬币的Poisson极限说明了这一点。

在独立同分布、方差有限且为正的条件下，标准化部分和依分布趋于$\mathcal{N}(0,1)$。证明把和的特征函数写成$n$个相同因子的乘积：独立性给出乘积，中心化消去一次项，方差为一固定二次项，二阶矩有限经支配收敛使累计余项可以忽略，连续性定理把特征函数的收敛变回分布的收敛。各项分布不同时，以行内总方差$s_n^2$归一化，Lindeberg条件要求大值贡献的方差份额趋于零；Lyapunov条件用$2+\delta$阶矩换来更方便的核验。这些都是充分条件，各项为高斯而最大项占一半方差的例子说明它们不是必要的。

随机向量的情形由方向投影化为标量：正方差方向用标量定理，零方差方向的投影几乎必然为零，再由Cramér--Wold判据合成协方差为$\bm{\Sigma}$的联合高斯极限，$\bm{\Sigma}$可以奇异。骰子频率的极限就是一个集中在超平面上的退化高斯分布。

\begin{BookInlineTable}
 \centering
 \caption{本章的高斯极限与误差界}
 \label{tab:clt-summary}
 \begin{tabular}{lll}
 \toprule
 结论 & 对象 & 条件 \\
 \midrule
 定理~\ref{thm:clt-iid} & $Z_n$ & 独立同分布，$0<\sigma^2<\infty$ \\
 定理~\ref{thm:clt-lindeberg} & $S_n/s_n$ & 行内独立，Lindeberg条件 \\
 定理~\ref{thm:clt-lyapunov} & $S_n/s_n$ & 行内独立，$\Lambda_n(\delta)\to0$ \\
 定理~\ref{thm:clt-multivariate} & $\sqrt{n}(\bar{\bm{X}}_n-\bm{\mu})$ & 独立同分布，二阶矩有限 \\
 定理~\ref{thm:clt-berry-esseen} & $D_n\leq C\rho/(\sigma^3\sqrt{n})$ & 独立同分布，$\rho<\infty$ \\
 \bottomrule
 \end{tabular}
\end{BookInlineTable}

极限的存在不等于有限样本的精度。Berry--Esseen界在三阶绝对矩有限时给出对所有阈值一致的绝对误差$C\rho/(\sigma^3\sqrt{n})$，速度$1/\sqrt{n}$在格点分布上不能改进；对光滑的分布它可能偏保守。它不控制小尾概率的相对误差，偏斜分布在三四个标准差以外的尾概率可以与高斯近似相差数十倍；三阶矩无穷时，定理仍然成立，误差界却无从调用。

使用这些结论时，先确认中心与尺度，再看独立结构是同分布、三角阵还是向量，然后按要回答的问题核对矩条件：只问极限，二阶矩或Lindeberg条件就够了；要给定样本量下的误差，需要三阶矩；要罕见事件的概率，应回到精确分布或指数矩上界。第~\ref{chap:monte-carlo-integration-error}章将用本章的结论描述独立Monte Carlo积分的误差。
